\subsubsection{Nonlinear generation using Shao--Wu moment contraction}
\label{sec:implemented-shao-wu}

The desired result is a decreasing influence of the initial history. Imagine
running the same equation twice, starting with different past observations but
using exactly the same future noise values. We want the expected squared
distance between the two trajectories to approach zero at a controlled rate.
For a mechanism with $L$ lags and $N$ channels, the initial history contains
the last $L$ observations of every channel. Sharing the future noise isolates
the effect of this history; independently simulated trajectories also differ
because of their fresh innovations.

\paragraph{Result sought: geometric forgetting of the past.}
\label{sec:implemented-shao-wu-theory}

\citet[Section~3, Eq.~(3.1), electronic reprint p.~5]{shao2007spectral}
formalize this idea as \emph{geometric moment contraction} (GMC). In their
comparison, two stationary versions have independently drawn innovation
histories through the origin $\tau$ and share all innovations after it.
For the second-moment version used here, their expected squared distance obeys
\begin{equation}
 \mathbb E\|\mathbf x_{\tau+v}-\mathbf x'_{\tau+v}\|_2^2
 \leq C_{\mathrm{GMC}}\rho_{\mathrm{GMC}}^{v},
 \qquad C_{\mathrm{GMC}}>0,\quad 0<\rho_{\mathrm{GMC}}<1.
 \label{eq:implemented-gmc-forgetting}
\end{equation}
The expectation averages the squared distance over repeated realizations of
this comparison. Since $\rho_{\mathrm{GMC}}<1$, the upper bound tends to zero.
Individual steps can fluctuate while remaining under a decaying bound in
expectation. Recent lags still influence each new observation, and innovations
sustain variation. Theorem~5.1 gives sufficient conditions for this property
and for a causal strictly stationary solution: a process determined by present
and past innovations whose joint distributions are unchanged by time shifts.
Our initialization procedure below uses the same contraction condition to
control the effect of a finite starting history.

\paragraph{Conditions supplied by Theorem 5.1.}
Let $\mathbf{s}_t$ stack the $L$ most recent vectors of $N$ series and write
\begin{equation}
 \mathbf x_{t+1}=\mathcal R_{\mathcal T}(\mathbf{s}_t;
                                      \boldsymbol\varepsilon_{t+1}).
 \label{eq:implemented-gmc-recurrence}
\end{equation}
We use the following second-moment specialization of
\citet[Section~5, Theorem~5.1 and its proof, electronic reprint
pp.~11--12]{shao2007spectral}:
\begin{enumerate}
\item \emph{A fixed update rule and independent innovations.}
The function $\mathcal R_{\mathcal T}$ is measurable, so its probabilistic
quantities are well-defined, and the same rule applies at every time step.
The innovations are independent and identically distributed over time.
\item \emph{A finite second moment at a reference history.}
At least one fixed history must give an output with finite expected squared
magnitude. Our construction checks the zero history, requiring
$\|\mathcal R_{\mathcal T}(\mathbf 0;\boldsymbol\varepsilon)\|_{L^2}<\infty$,
where $\|Z\|_{L^2}=(\mathbb E\|Z\|_2^2)^{1/2}$.
\item \emph{Limited sensitivity to the past.}
For every pair of histories, the difference between their next outputs under
the same innovation must satisfy
\begin{equation}
\begin{split}
 \|\mathcal R_{\mathcal T}(\mathbf{s};\boldsymbol\varepsilon)
      -\mathcal R_{\mathcal T}(\mathbf{s}';\boldsymbol\varepsilon)\|_{L^2}
 &\leq\sum_{d=0}^{L-1}a_d^{\mathrm{GMC}}
                   \|\mathbf{s}_d-\mathbf{s}'_d\|_2,\\
 a_d^{\mathrm{GMC}}&\geq0,\qquad
 c_{\mathrm{GMC}}:=\sum_{d=0}^{L-1}a_d^{\mathrm{GMC}}<1.
\end{split}
\label{eq:implemented-gmc-condition}
\end{equation}
Here $\mathbf{s}_d\in\mathbb R^N$ contains all channels at one lag.
Each nonnegative $a_d^{\mathrm{GMC}}$ bounds how much a change at that lag
can affect the next output. Their sum below one controls the combined feedback
as the history is repeatedly updated.
\end{enumerate}
These conditions apply globally, across all possible lag values. The full
theorem permits other moment orders; the implementation uses order two.

\paragraph{Our construction: making the hypotheses computable.}
\label{sec:implemented-shao-wu-construction}

The following steps, implemented in
\path{src/tsfmxai/generators/shao_wu.py}, turn the theorem's conditions into
checks on a sampled equation. The random tree sampler defines which mechanisms
we draw. Each mechanism has a \emph{verification record}, named
\texttt{ShaoWuCertificate} in the code. It stores the calculated lag
sensitivities, their sum, the amplitude and noise bounds, and the initialization
error bound. This is
the meaning of ``certificate'' here: recorded quantities used to check the
stated sufficient conditions and the chosen initialization tolerance.

\subparagraph{1. Bound the lag terms and keep the same rule over time.}
For each target, sample a base tree with parameter and lag leaves, addition,
and multiplication. Replace each lag leaf $x_{j,t-\ell}$ with
$\tanh(x_{j,t-\ell})$. For every real input, $\tanh$ lies in $(-1,1)$ and
its derivative has magnitude at most one. Consequently, we can bound both
the size of each transformed lag and its sensitivity before combining the
terms through sums and products. With finite coefficients, this gives finite
bounds for the complete deterministic part, also called the drift.

The tree uses fixed parameters and lagged observations, keeping the same
measurable update rule at every step. We add i.i.d. Gaussian innovations with
covariance $\sigma_\varepsilon^2 I$. This supplies condition~1; the finite
drift and noise moments will supply condition~2. The derivative bounds prepare
condition~3, whose strict feedback limit is enforced in step~3.

\subparagraph{2. Calculate size and sensitivity bounds from the leaves upwards.}
For each subexpression $e$, the recursive calculation returns an amplitude
bound $A(e)$ and a matrix $D(e)$ of nonnegative absolute derivative bounds,
indexed by source series and lag offset. Table~\ref{tab:ast-global-bounds}
gives the implemented rules. In practical terms, $A(e)$ answers how large the
deterministic expression can become, while $D(e)$ bounds its response to a
change in each past input. We calculate these quantities recursively until
we reach the root of each output tree. Their global validity allows us to
check the theorem's conditions before simulating a trajectory.

\begin{table}[H]
\centering\small
\begin{tabular}{@{}lll@{}}
\toprule
Subexpression & Amplitude bound $A$ & Derivative bounds $D$ \\
\midrule
Constant $c$ & $|c|$ & $0$ \\
Lag leaf & $+\infty$ & $1$ at its source/lag, $0$ elsewhere \\
$\tanh(e)$ & $1$ & $D(e)$ \\
$e_1+e_2$ & $A(e_1)+A(e_2)$ & $D(e_1)+D(e_2)$ \\
$e_1e_2$ & $A(e_1)A(e_2)$ & $A(e_2)D(e_1)+A(e_1)D(e_2)$ \\
\bottomrule
\end{tabular}
\caption{Conservative global bound propagation in our Shao--Wu construction.
The automatic sampler wraps every lag leaf in $\tanh$ before products are
bounded. A zero derivative contributes zero even if the other amplitude
bound is infinite.}
\label{tab:ast-global-bounds}
\end{table}

For a sum, adding the two bounds covers the largest possible combined effect.
For a product, the product rule explains the two derivative terms: a change
in one factor is multiplied by the size of the other. This is why bounding
the lag values before forming products is useful.

\subparagraph{3. Reduce the combined sensitivity until its bound is below one.}
For lag offset $d$, stack the bounds from all targets into a nonnegative
$N\times N$ matrix $\mathbf D_d$, whose entry $(k,j)$ bounds the effect of
source $j$ on target $k$. The norm $\|\mathbf D_d\|_2$, its largest singular
value, bounds how much changing the whole lag vector can change the output
vector. It accounts for feedback between channels. Define
\begin{equation}
 \widetilde a_d=\|\mathbf D_d\|_2,\qquad
 \widetilde c=\sum_{d=0}^{L-1}\widetilde a_d,
 \qquad
 \eta_{\mathrm{GMC}}=
 \eta_{\mathrm{safe}}\min\{1,g_\star/\widetilde c\}.
 \label{eq:implemented-gmc-scaling}
\end{equation}
Here $\widetilde c$ adds the sensitivity bounds over all lags. If
$\widetilde c=0$, the
minimum in Eq.~\eqref{eq:implemented-gmc-scaling} is defined as 1. The
defaults are $g_\star=0.8$ and $\eta_{\mathrm{safe}}=0.95$.
Multiply every output tree by $\eta_{\mathrm{GMC}}$, giving
$a_d^{\mathrm{GMC}}=\eta_{\mathrm{GMC}}\widetilde a_d$ and
$c_{\mathrm{GMC}}\leq0.76<1$.

This multiplication reduces every lag sensitivity by the same factor. The
resulting total directly meets the strict inequality in condition~3.
To connect derivatives to that condition, consider the line joining two
histories: the global derivative bounds limit the change along the entire
line. Adding the lagwise bounds gives
Eq.~\eqref{eq:implemented-gmc-condition}. The additive noise cancels because
the comparison uses exactly the same innovation on both sides.

\subparagraph{4. Check the output moment and choose the initialization length.}
Let $B_{\mathrm{GMC}}$ be the Euclidean norm of the final componentwise
amplitude bounds. The stationary solution satisfies
$\|\mathbf x_t\|_{L^2}\leq B_{\mathrm{GMC}}+\sigma_\varepsilon\sqrt N$.
The drift is bounded and Gaussian noise has finite variance, so the next
output at the zero history also has a finite second moment. This verifies
condition~2 and completes the checks needed to apply the theorem.

We then use the contraction to decide how many preliminary steps to simulate,
a procedure called \emph{burn-in}. Starting each lag block in
$[-r_0,r_0]^N$ gives the following bound on its difference from a stationary
version driven by the same future innovations:
\begin{equation}
\begin{gathered}
 M_0=r_0\sqrt N+B_{\mathrm{GMC}}+\sigma_\varepsilon\sqrt N,\\
 \max_{0\leq d<L}
 \|\mathbf x_{b_{\mathrm{burn}}-d}-\mathbf x^{\mathrm{stat}}_{b_{\mathrm{burn}}-d}\|_{L^2}
 \leq M_0c_{\mathrm{GMC}}^{K}.
\end{gathered}
 \label{eq:implemented-gmc-burnin}
\end{equation}
The first term in $M_0$ bounds our chosen initial values; the remaining terms
bound a stationary observation. Their sum bounds the initial difference.
Here $b_{\mathrm{burn}}=L-1+KL$ counts burn-in steps from an initial history ending at zero.
The maximum lagwise $L^2$ error contracts by at most $c_{\mathrm{GMC}}$
after each complete block of $L$ steps: once a block is generated, every lag
entering the next block already has the reduced error bound. The additional
$L-1$ steps conservatively cover every observation in the retained history.
Thus more completed blocks make the starting history less influential.
The bound concerns each lag block, with
the maximum outside the expectation; a bound on the entire stacked
Euclidean state gains a factor $\sqrt L$.

For $0<c_{\mathrm{GMC}}<1$ and $M_0>\delta_{\mathrm{init}}$, choose
$K=\lceil\log(\delta_{\mathrm{init}}/M_0)/\log(c_{\mathrm{GMC}})\rceil$.
The code uses $K=0$ when the initial bound already meets the tolerance, and
$K=1$ when the gain is zero and a contraction block is needed. The default
per-lag tolerance is $\delta_{\mathrm{init}}=10^{-4}$. Initialization simulates
the Gaussian-driven recurrence for $b_{\mathrm{burn}}$ steps and retains the final $L$
observations. Fresh, independently seeded innovations generate the returned
trajectory. This gives approximate stationary initialization with a stated
remaining error. The theorem establishes the stationary solution; the burn-in
calculation specifies how closely our finite initialization approaches it.

\subparagraph{A worked construction.}
Consider our example drift $0.3\tanh(x_t)\tanh(x_{t-1})$.
Its amplitude bound is 0.3, and both lag derivative bounds are 0.3, giving
$\widetilde c=0.6$. The default safety factor yields the executed recurrence
\begin{equation}
 x_{t+1}=0.285\tanh(x_t)\tanh(x_{t-1})+\varepsilon_{t+1},
 \qquad \varepsilon_{t+1}\sim\mathcal N(0,0.05^2).
 \label{eq:implemented-gmc-example}
\end{equation}
Now $c_{\mathrm{GMC}}=0.57$, $M_0=1.335$ for $r_0=1$, and $K=17$ gives
$b_{\mathrm{burn}}=35$ burn-in steps. The gain 0.57 satisfies condition~3;
the bounded drift and Gaussian variance satisfy condition~2. The 35 preliminary
steps bring the bound for every retained lag below $10^{-4}$. The verification
record saves these values alongside the scaled AST, making the reasoning
inspectable for each generated mechanism. Tests check the analytic example,
reproducibility, and consistency with sampled numerical derivatives.

\subparagraph{Scope of this construction.}
The implemented constructor requires both finite global drift amplitudes
and finite global derivative bounds. This selects a convenient bounded-drift
subfamily of the theorem: unbounded polynomial feedback is rejected, while
unbounded linear drift can be handled by the separate VAR family. Conservative
bounds and scaling can weaken useful signal relative to the innovations.
Accordingly, low-amplitude diagnostics remain relevant when the GMC conditions
hold. Gaussian innovations can still produce observations beyond a finite
amplitude threshold. The guarantee concerns the influence of the past and
the process's moments, with fresh noise continuing to drive its fluctuations.

The inequalities above describe the construction in exact arithmetic. The
implementation evaluates its bounds and singular values in floating point,
retains a strict gain margin, and checks the recorded bounds and sampled
derivatives. These checks provide numerical validation of the implementation;
their precision is governed by the numerical routines used.

\paragraph{Implications of independent innovations.}
\label{sec:implemented-shao-wu-independent-noise}
The shared-noise comparison uses the same sampled innovation value in both
trajectories at each future step, isolating the influence of their different
pasts \citep[Section~3, Eq.~(3.1), electronic reprint p.~5]{shao2007spectral}.
For dataset generation, each trajectory can instead receive a separate
innovation sequence. We consider the same fixed mechanism and innovation law
in both runs, with innovations i.i.d. over time, independent between runs, and
independent of their initial histories.

The relevant guarantee then concerns their probability distributions.
\citet[Section~2, Theorem~2, p.~427]{wu2004iterated} establish convergence of
backward random-function iterations under a finite-moment condition and
geometric contraction. Their explanation on p.~426 transfers this result to
convergence in distribution of forward simulations, and Remark~2 on p.~428
constructs the stationary process. For our finite-lag system, the blockwise
contraction used above controls the augmented history and yields this
forgetting of initialization. Applying the result separately to two runs
gives the following consequence:
\begin{equation}
 \operatorname{Law}(\mathbf x_{\tau+v})\Rightarrow\pi_{\mathrm{GMC}},
 \qquad
 \operatorname{Law}(\mathbf x'_{\tau+v})\Rightarrow\pi_{\mathrm{GMC}},
 \qquad v\to\infty.
 \label{eq:implemented-gmc-independent-laws}
\end{equation}
Here $\pi_{\mathrm{GMC}}$ is the stationary distribution of the observation
vector for that mechanism. Each run follows the same transition rule and
noise law, so both approach the same statistical regime even with independent
noise realizations. The contraction and moment assumptions remain essential.

Independent trajectories can still have different values at every time.
For two copies independently initialized in stationarity with finite second
moments, expanding the squared distance gives
\begin{equation}
 \mathbb E\|\mathbf x_t-\mathbf x'_t\|_2^2
 =2\sum_{k=1}^{N}\operatorname{Var}_{\pi_{\mathrm{GMC}}}(x_{k,t}).
 \label{eq:implemented-gmc-independent-distance}
\end{equation}
Independence makes the cross-products factor into products of means, leaving
twice the sum of the channel variances. This is an elementary consequence of
independent stationary sampling, distinguished here from the cited convergence
theorem. Positive stationary variance therefore sustains a positive expected
squared distance while both copies have the same distribution.

In our implementation, \texttt{sample\_shao\_wu\_trajectories} assigns each
trajectory separate initialization and future-noise streams derived from the
root seed and trajectory index. Burn-in controls each trajectory's remaining
initialization effect through Eq.~\eqref{eq:implemented-gmc-burnin}. Shared-noise
pairs serve to study sensitivity to the past; independent runs provide varied,
reproducible realizations of one mechanism. A persistent distance between
independent runs is compatible with this guarantee. Distributional comparisons
should use repeated runs, while the explosion and tail diagnostics continue
to assess individual finite trajectories.
